% Reference-tune comparison: the best GENIE configuration
%   N/LFG + F_A^LQCD + hA2018  (= LFG26_24a_00_000)
% against the AR23 and MicroBooNE reference tunes; 6-panel figure followed by its
% chi2/ndf table. Per-panel PDFs are co-located as reference_tune_{a..f}.pdf;
% prepend your paper's path (e.g. plots/) to the \includegraphics keys.
% chi2/ndf (p-value) rows are reused verbatim from chisq_table.tex (tab:chisq).
% Requires: amsmath (for \text), graphicx, overpic, booktabs.
% Panels: (a) dpT (dalphaT<45 deg), (b) dpT (135<dalphaT<180), (c) dpT,
%         (d) dalphaT, (e) p_mu, (f) p_p.

% ============== best config vs AR23 / MicroBooNE reference tunes ==============
\begin{figure*}[htbp]
\centering
\begin{minipage}{0.3\textwidth}
\begin{overpic}[width=\linewidth]{reference_tune_a.pdf}
\put(24,56){\footnotesize\text{(a)}}
\end{overpic}
\end{minipage}
%\hfill
\begin{minipage}{0.3\textwidth}
\begin{overpic}[width=\linewidth]{reference_tune_b.pdf}
\put(24,56){\footnotesize\text{(b)}}
\end{overpic}
\end{minipage}
%\hfill
\begin{minipage}{0.3\textwidth}
\begin{overpic}[width=\linewidth]{reference_tune_c.pdf}
\put(24,56){\footnotesize\text{(c)}}
\end{overpic}
\end{minipage}

%\hfill
\vspace{0.5cm}
\begin{minipage}{0.3\textwidth}
\begin{overpic}[width=\linewidth]{reference_tune_d.pdf}
\put(24,56){\footnotesize\text{(d)}}
\end{overpic}
\end{minipage}
\begin{minipage}{0.3\textwidth}
\begin{overpic}[width=\linewidth]{reference_tune_e.pdf}
\put(24,56){\footnotesize\text{(e)}}
\end{overpic}
\end{minipage}
%\hfill
\begin{minipage}{0.3\textwidth}
\begin{overpic}[width=\linewidth]{reference_tune_f.pdf}
\put(24,56){\footnotesize\text{(f)}}
\end{overpic}
\end{minipage}
\caption{The fully theory-driven GENIE configuration N/LFG $+\ F_A^{\rm LQCD}\ +$ hA2018 (blue) compared with the AR23 (orange) and MicroBooNE (green) reference tunes on the MicroBooNE CC1$\mu$1p / CC1$\mu$Np argon cross sections (black points). Panels: (a) $\delta p_T$ for $\delta\alpha_T<45^\circ$ and (b) for $135^\circ<\delta\alpha_T<180^\circ$ (two slices of the 2D $\delta p_T$--$\delta\alpha_T$ measurement); (c) $\delta p_T$; (d) $\delta\alpha_T$; (e) $p_\mu$; (f) $p_p$. The reference tunes are drawn dashed in the projected panels (c)--(f). The N/LFG $+\ F_A^{\rm LQCD}\ +$ hA2018 prediction gives the lowest $\chi^2/{\rm ndf}$ of the three in every observable (see Table~\ref{tab:chisq_reference_tune}).}
\label{fig:reference_tune}
\end{figure*}

\begin{table}[htbp]
\caption{$\chi^2/{\rm ndf}$ ($p$-value) relative to MicroBooNE data for Fig.~\ref{fig:reference_tune}.}
\label{tab:chisq_reference_tune}
\centering
\resizebox{\columnwidth}{!}{
\begin{tabular}{lrrrrr}
\toprule
Configuration & $\delta p_T$ vs. $\delta \alpha_T$ & $\delta p_T$ & $\delta \alpha_T$ & $p_{\mu}$ & $p_p$ \\
\midrule
N/LFG + $F_A^{\rm LQCD}$ + hA2018 & 35.98/49  (0.92) & 6.78/13  (0.91) & 3.47/7  (0.84) & 16.29/26  (0.93) & 16.05/15  (0.38)\\
AR23                              & 38.05/49  (0.87) & 10.63/13  (0.64) & 7.56/7  (0.37) & 28.96/26  (0.31) & 24.07/15  (0.06)\\
MicroBooNE Tune                   & 42.61/49  (0.73) & 7.31/13  (0.89) & 3.96/7  (0.78) & 24.59/26  (0.54) & 17.38/15  (0.30)\\
\bottomrule
\end{tabular}}
\end{table}
